\chapter{KAM Proof Architecture}
\label{app:kam-proof-architecture}

This appendix gives a more explicit view of the proof mechanism behind the KAM
statement in Chapter~\ref{chap:perturbation-chaos}.  It is not a substitute for
a full monograph proof, but it records the analytic estimates, the role of the
twist condition, and the reason that the theorem is a Newton method on
Hamiltonians rather than an ordinary first-order perturbation calculation.

\section{Local Analytic Model}
\label{sec:kam-local-analytic-model}

Let \(n\) be the number of degrees of freedom.  The unperturbed phase space is
\[
        \mathcal{D}\times \mathbb{T}^n,\qquad
        (I,\theta)=(I_1,\dots,I_n,\theta^1,\dots,\theta^n),
\]
where \(\mathcal{D}\subset \R^n\) is an open action domain and
\(\mathbb{T}^n=\R^n/(2\pi\Z)^n\).  The symplectic form is
\[
        \om=\sum_{a=1}^n \dd I_a\wedge \dd \theta^a,
\]
and the Hamiltonian is
\[
        \Ham(I,\theta)=h(I)+\eps f(I,\theta).
\]
Here \(h\) is the integrable Hamiltonian, \(f\) is a perturbation, and \(\eps\)
is a small dimensionless parameter.  The frequency map of the unperturbed
motion is
\[
        \omega(I)=\nabla_I h(I).
\]
If \(\eps=0\), the torus \(I=I_0\) is invariant and the angle motion is
\[
        \theta(t)=\theta(0)+\omega(I_0)t .
\]

\begin{assumption}[Analyticity and complex strips]
For the proof sketch in this appendix, assume that \(h\) and \(f\) extend
holomorphically to a complex neighborhood
\[
        \mathcal{D}_\rho\times \mathbb{T}^n_\sigma
        =
        \{I\in\mathbb{C}^n:\operatorname{dist}(I,\mathcal{D})<\rho\}
        \times
        \{\theta\in\mathbb{C}^n/(2\pi\Z)^n:|\operatorname{Im}\theta|<\sigma\}.
\]
The constants \(\rho\) and \(\sigma\) are analytic radii.  Estimates lose part
of these radii at each coordinate change; that loss is the analytic price paid
for controlling derivatives and Fourier tails.
\end{assumption}

For an analytic function
\[
        F(I,\theta)=\sum_{k\in\Z^n}F_k(I)e^{i k\cdot\theta},
\]
a convenient norm is
\[
        \norm{F}_{\rho,\sigma}
        =
        \sum_{k\in\Z^n}\sup_{I\in\mathcal{D}_\rho}|F_k(I)|e^{|k|\sigma},
        \qquad |k|=|k_1|+\cdots+|k_n|.
\]
The exponential weight records the analytic decay of Fourier coefficients.
The KAM proof repeatedly trades this decay against small denominators
\(k\cdot\omega\).

\section{The Homological Equation}
\label{sec:kam-homological-equation}

The first calculation is the one canonical perturbation theory already
suggests.  Fix an action \(I_0\) and write
\[
        N(I)=h(I),\qquad P(I,\theta)=\eps f(I,\theta).
\]
We seek a near-identity canonical transformation generated by a Hamiltonian
\(\chi(I,\theta)\).  If \(\Phi_\chi^1\) denotes the time-one map of the
Hamiltonian flow of \(\chi\), then Lie expansion gives
\[
        \Ham\circ \Phi_\chi^1
        =
        \Ham+\Poisson{\Ham}{\chi}
        +\frac{1}{2}\Poisson{\Poisson{\Ham}{\chi}}{\chi}+\cdots .
\]
To first order in the perturbation, one tries to choose \(\chi\) so that the
angle-dependent part of \(P\) is removed:
\[
        \Poisson{N}{\chi}+P-\langle P\rangle=0,
        \qquad
        \langle P\rangle(I)=\frac{1}{(2\pi)^n}\int_{\mathbb{T}^n}P(I,\theta)\,\dd^n\theta .
\]
Since \(N=h(I)\),
\[
        \Poisson{N}{\chi}
        =
        \nabla_I h(I)\cdot \pd_\theta\chi
        =
        \omega(I)\cdot \pd_\theta\chi .
\]
At \(I=I_0\), the Fourier coefficients of the equation are
\[
        i(k\cdot \omega_0)\chi_k(I_0)+P_k(I_0)=0,
        \qquad k\ne0,\qquad \omega_0=\omega(I_0),
\]
so that
\[
        \chi_k(I_0)=\frac{iP_k(I_0)}{k\cdot\omega_0}.
\]
This formula is the source of the small-divisor problem.  The numerator
\(P_k\) decays exponentially in \(|k|\) if the perturbation is analytic, but
the denominator can be very small when \(k\cdot\omega_0\) nearly vanishes.

\begin{definition}[Diophantine frequency]
A vector \(\omega\in\R^n\) is Diophantine with constants
\(\gamma>0\) and \(\tau>n-1\) if
\[
        |k\cdot\omega|\ge \frac{\gamma}{|k|^\tau}
        \qquad \text{for all }0\ne k\in\Z^n .
\]
The set of such vectors is denoted \(\mathrm{DC}(\gamma,\tau)\).
\end{definition}

If \(\omega_0\in\mathrm{DC}(\gamma,\tau)\), then
\[
        |\chi_k(I_0)|
        \le
        \gamma^{-1}|k|^\tau |P_k(I_0)|.
\]
The factor \(|k|^\tau\) is polynomial.  The analytic decay of \(P_k\) can
absorb it if one gives up a smaller angle strip.  More explicitly, for
\(0<\sigma'<\sigma\),
\[
        |k|^\tau e^{-|k|\sigma}
        =
        |k|^\tau e^{-|k|(\sigma-\sigma')}e^{-|k|\sigma'}
        \le
        C_{\tau,n}(\sigma-\sigma')^{-(\tau+n)}e^{-|k|\sigma'} .
\]
Thus the generator satisfies an estimate of the form
\[
        \norm{\chi}_{\rho,\sigma'}
        \le
        C\gamma^{-1}(\sigma-\sigma')^{-(\tau+n)}
        \norm{P}_{\rho,\sigma}.
\]
The constants are not the conceptual point.  The point is the structure:
small divisors are controlled by Diophantine arithmetic, and derivatives are
controlled by sacrificing analytic width.

\section{Twist, Frequencies, And What Is Preserved}
\label{sec:kam-twist-fixed-frequency}

The KAM theorem does not say that every original torus \(I=I_0\) remains
exactly at the same action.  A perturbation changes the average Hamiltonian,
and the frequency map moves.  To track a torus with a prescribed frequency
\(\omega_*\), one needs a nondegenerate relation between actions and
frequencies.

\begin{assumption}[Kolmogorov nondegeneracy]
On the action domain under consideration,
\[
        \det \pd_I^2 h(I)\ne0 .
\]
Equivalently, the frequency map \(I\mapsto \omega(I)\) is locally invertible.
\end{assumption}

After the first averaging step the new integrable part is
\[
        h_+(I)=h(I)+\langle P\rangle(I)+\text{higher order terms}.
\]
The action whose frequency equals \(\omega_*\) generally shifts from \(I_0\) to
\(I_0+\Delta I\).  Linearizing,
\[
        \omega_*=
        \nabla h_+(I_0+\Delta I)
        =
        \nabla h(I_0)
        +
        \pd_I^2h(I_0)\Delta I
        +
        \nabla\langle P\rangle(I_0)
        +\cdots .
\]
If \(\nabla h(I_0)=\omega_*\), then
\[
        \Delta I
        =
        -\bigl(\pd_I^2h(I_0)\bigr)^{-1}\nabla\langle P\rangle(I_0)
        +\cdots .
\]
The twist condition is the assumption that this adjustment can be solved for
uniquely.  Without it, KAM theorems still exist in modified forms, but some
other nondegeneracy condition must replace it.

\section{Isoenergetic And Map Versions}
\label{sec:kam-isoenergetic-map-versions}

The nondegeneracy condition depends on which object is being preserved.  The
Kolmogorov theorem fixes a frequency vector and constructs a nearby invariant
torus with that frequency.  In applications one often fixes an energy surface
instead.  Then frequencies are meaningful only up to a common rescaling of
time, because the orbit on a torus with frequency \(\omega\) and the orbit with
frequency \(c\omega\), \(c>0\), have the same unparametrized winding direction.

\begin{definition}[Isoenergetic nondegeneracy]
The integrable Hamiltonian \(h(I)\) satisfies the isoenergetic nondegeneracy
condition at \(I\) if
\[
        \det
        \begin{pmatrix}
        \pd_I^2h(I) & \nabla h(I)\\
        \nabla h(I)^T & 0
        \end{pmatrix}
        \ne 0 .
\]
\end{definition}

This condition says that, after restricting to an energy surface, the
frequency direction still changes nontrivially with the actions.  It is the
natural condition when the geometry of a fixed energy hypersurface is the main
object.

For a two-dimensional area-preserving twist map with coordinates \((\theta,I)\)
and unperturbed form
\[
        I_{m+1}=I_m,\qquad
        \theta_{m+1}=\theta_m+\Omega(I_m),
\]
the corresponding nondegeneracy condition is
\[
        \frac{\dd \Omega}{\dd I}\ne0 .
\]
This is the map version of twist.  It is what makes the standard map a useful
laboratory: for \(K=0\),
\[
        p_{m+1}=p_m,\qquad q_{m+1}=q_m+p_m ,
\]
so the rotation number changes with \(p\).  KAM curves of the perturbed map are
organized by their irrational rotation numbers.

\begin{center}
\begin{tabular}{p{0.25\textwidth} p{0.34\textwidth} p{0.29\textwidth}}
\hline
setting & preserved object & nondegeneracy condition\\
\hline
Hamiltonian flow, fixed frequency & torus with frequency \(\omega\) &
\(\det \pd_I^2h\ne0\)\\
Hamiltonian flow, fixed energy & frequency direction on an energy surface &
bordered determinant above is nonzero\\
area-preserving twist map & invariant curve with rotation number \(\rho\) &
\(\dd\Omega/\dd I\ne0\)\\
\hline
\end{tabular}
\end{center}

\section{One KAM Step}
\label{sec:one-kam-step}

The proof is organized as an iteration.  At the \(m\)-th stage the Hamiltonian
has the form
\[
        \Ham_m(I,\theta)
        =
        h_m(I)+P_m(I,\theta)
\]
on a complex domain
\(\mathcal{D}_{\rho_m}\times\mathbb{T}^n_{\sigma_m}\), with
\(\norm{P_m}\) small.  The step has four parts.

\begin{enumerate}
\item \emph{Truncate Fourier modes.}  Write \(P_m\) as low modes
\(|k|\le K_m\) plus a high-mode tail.  Analyticity makes the tail
exponentially small:
\[
        \norm{P_m^{>|K_m|}}_{\sigma_m-\delta_m}
        \le
        e^{-K_m\delta_m}\norm{P_m}_{\sigma_m}.
\]
\item \emph{Solve the homological equation.}  Remove the low-mode
angle-dependent part by solving
\[
        \omega_m(I_*)\cdot\pd_\theta\chi_m
        +
        P_m^{\le K_m}(I_*,\theta)
        -
        \langle P_m^{\le K_m}\rangle(I_*)
        =0 .
\]
\item \emph{Apply the canonical transformation.}  Use the time-one flow
\(\Phi_{\chi_m}^1\).  The transformed Hamiltonian is
\[
        \Ham_{m+1}=\Ham_m\circ\Phi_{\chi_m}^1=h_{m+1}+P_{m+1}.
\]
\item \emph{Recentre the action.}  Use the twist condition to keep the
frequency of the selected torus equal to the desired Diophantine vector.
\end{enumerate}

The new perturbation contains a quadratic part, a Fourier-tail part, and
domain-loss errors:
\[
        \norm{P_{m+1}}
        \lesssim
        C\frac{\norm{P_m}^2}{\gamma^2\delta_m^{2(\tau+n)+1}}
        +
        C e^{-K_m\delta_m}\norm{P_m}.
\]
The parameters \(K_m\) and \(\delta_m\) are chosen so that the right side is
much smaller than \(\norm{P_m}\).  This is the Newton character of the method:
after the resonant obstructions are excluded and the domain losses are budgeted,
the perturbation decreases roughly quadratically.

\begin{figure}[htbp]
\centering
\begin{tikzpicture}[
    box/.style={draw, rounded corners=2pt, minimum width=2.7cm, minimum height=0.9cm, align=center},
    arrow/.style={-{Stealth[length=2.5mm]}, thick}
]
\node[box] (h0) at (0,0) {\(\Ham_m=h_m+P_m\)};
\node[box] (trunc) at (3.6,0) {truncate\\ Fourier modes};
\node[box] (hom) at (7.2,0) {solve\\ homological equation};
\node[box] (map) at (3.6,-1.8) {canonical\\ time-one map};
\node[box] (new) at (7.2,-1.8) {\(\Ham_{m+1}=h_{m+1}+P_{m+1}\)};
\node[box] (shift) at (0,-1.8) {recentre\\ fixed frequency};
\draw[arrow] (h0) -- (trunc);
\draw[arrow] (trunc) -- (hom);
\draw[arrow] (hom) -- (map);
\draw[arrow] (map) -- (new);
\draw[arrow] (new) -- (shift);
\draw[arrow] (shift) -- (h0);
\end{tikzpicture}
\caption{A KAM step is a controlled averaging step followed by a recentring of
the action.  The loop contracts because the remaining perturbation is
quadratic up to an exponentially small Fourier tail.}
\label{fig:kam-newton-step}
\end{figure}

\section{Cantor Set And Measure Estimate}
\label{sec:kam-cantor-measure}

The theorem keeps only those frequencies that remain safely away from
resonance.  For a bounded frequency domain \(\Omega\subset\R^n\), define
\[
        \Omega_{\gamma,\tau}
        =
        \left\{\omega\in\Omega:
        |k\cdot\omega|\ge \frac{\gamma}{|k|^\tau}
        \text{ for every }0\ne k\in\Z^n\right\}.
\]
The complement is a union of resonance strips
\[
        \mathcal{R}_k
        =
        \left\{\omega\in\Omega:
        |k\cdot\omega|<\frac{\gamma}{|k|^\tau}\right\}.
\]
For a fixed \(k\), this strip has thickness of order
\(\gamma |k|^{-(\tau+1)}\) in the direction normal to the hyperplane
\(k\cdot\omega=0\).  Summing over \(k\ne0\),
\[
        \operatorname{meas}(\Omega\setminus\Omega_{\gamma,\tau})
        \le
        C_\Omega \gamma
        \sum_{0\ne k\in\Z^n}|k|^{-(\tau+1)} .
\]
The sum converges when \(\tau>n-1\).  Therefore the excluded set has measure
\(O(\gamma)\).  The surviving set is large in measure but is not open: it is a
Cantor-like set obtained by removing resonance neighborhoods at all integer
orders.

\section{Precise Form Of The Conclusion}
\label{sec:kam-conclusion}

One common form of the theorem can now be read without hiding the main
ingredients.

\begin{proposition}[KAM theorem, schematic analytic version]
Let
\[
        \Ham(I,\theta)=h(I)+\eps f(I,\theta)
\]
be real analytic on
\(\mathcal{D}_\rho\times\mathbb{T}^n_\sigma\).  Assume that
\(\det \pd_I^2 h(I)\ne0\) on \(\mathcal{D}\).  Fix
\(\tau>n-1\) and \(\gamma>0\).  If \(\eps\norm{f}_{\rho,\sigma}\) is
sufficiently small compared with a constant determined by \(h,\rho,\sigma\),
\(\gamma\), and \(\tau\), then for each frequency in a slightly smaller
Diophantine set there is an invariant Lagrangian torus of the perturbed
Hamiltonian.  On that torus the motion is analytically conjugate to the linear
flow \(\dot\theta=\omega\).  The union of destroyed tori in the corresponding
action region has measure tending to zero with \(\eps\).
\end{proposition}

The theorem is deliberately a persistence theorem, not a global solution of
the perturbed dynamics.  It says that sufficiently irrational invariant tori
survive.  It does not say that resonant tori survive, that all trajectories are
regular, or that transport is absent through the gaps between tori.  Those
gaps are where resonant islands, separatrix splitting, cantori, and chaotic
transport enter the main text.

\section{Checklist For Applying KAM}
\label{sec:kam-application-checklist}

When using KAM reasoning in examples, check the following points explicitly.

\begin{enumerate}
\item Identify the integrable Hamiltonian \(h(I)\), the perturbation
\(\eps f(I,\theta)\), and the action domain.
\item Verify the nondegeneracy condition being used.  In Kolmogorov's theorem
this is the twist condition \(\det \pd_I^2h\ne0\).
\item State the arithmetic condition on the frequency vector and the constants
\(\gamma,\tau\).
\item Explain which tori are excluded because they are resonant or nearly
resonant.
\item Distinguish persistence of a torus from confinement of nearby orbits.
In two-dimensional area-preserving maps, invariant curves are barriers.  In
higher-dimensional Hamiltonian flows, surviving tori need not divide an energy
surface.
\item Separate mathematical smallness from physical smallness.  A physically
small perturbation can still have important resonant effects if the orbit is
near a low-order resonance.
\end{enumerate}

This checklist is the bridge from the proof architecture to the standard map,
the asteroid-belt resonances, and the breakdown mechanisms discussed in
Chapter~\ref{chap:perturbation-chaos}.
